% Resumen ejecutivo (rúbrica 5 pts, junto con la introducción). Una página, sin siglas sin definir.
La Red Metropolitana de Movilidad (RED) de Santiago está reemplazando sus buses diésel por buses eléctricos. Un bus eléctrico tiene autonomía limitada y, cuando su batería no alcanza para la jornada, debe ir a cargar a uno de los cinco electroterminales, que suman \PlazasTotales{} plazas. Este proyecto decide, para un día laboral, un sábado y un domingo, cuántos buses usar, qué corridas cubre cada uno y cuándo, dónde y cuánto recarga, al menor costo diario, que considera USD~\CostoBus{} por bus, USD~\CostoKmUnit{} por kilómetro, la energía según la hora (entre USD~\TarifaMin{} y \TarifaMax{} por kWh) y USD~\CostoSesion{} por cada sesión de carga.

Con los horarios públicos se reconstruyeron \CorridasLaboral{} corridas de día laboral en \RutasLaboral{} rutas, \CorridasSabado{} de sábado y \CorridasDomingo{} de domingo. Como un modelo exacto solo resolvió casos de hasta \CorridasRutasDiez{} corridas, el plan se obtiene con un método híbrido. Cada ruta se asocia al electroterminal más cercano, se arma un plan inicial siguiendo el orden de salida de las corridas y después se mejora tomando grupos de cuatro buses al azar, en los que uno intenta repartir sus corridas entre los otros tres mientras un modelo entero pequeño decide sus recargas. Un auditor independiente revisa cada plan.

Frente a un caso base que planifica cada ruta por separado y también recarga, el plan cubre las mismas corridas del día laboral con \BusesPlanLaboral{} buses en lugar de \BusesBaseLaboral{} (\DeltaPctBusesLaboral\,\%) y cuesta USD~\CostoPlanLaboral{} al día (\DeltaPctCostoLaboral\,\%). El sábado la flota baja \BajaBusesSabado\,\% y el domingo \BajaBusesDomingo\,\%. En una semana el costo cae \BajaSemana\,\%, unos USD~\AhorroSemanaMillones{} millones, y cerca de tres de cada cuatro dólares ahorrados corresponden a buses. La solución es aproximada y queda a \BrechaCuartetos\,\% del óptimo donde este se pudo demostrar.

El ahorro se mantiene entre \FisAhorroMin{} y \FisAhorroMax\,\% con baterías más chicas o mayor consumo, y entre \CostosAhorroMin{} y \CostosAhorroMax\,\% si el costo por bus o por kilómetro cambia un 20\,\%. Con tiempos de viaje que varían un 10\,\%, el \AtrasadasPlanDiez\,\% de las corridas del plan sale con más de cinco minutos de atraso, contra el \AtrasadasBaseDiez\,\% del caso base. Y para repetir el día laboral habría que reponer \ReponerGwhLaboral{}~GWh durante la noche, lo que exigiría \ConexionesNecesariasLaboral{} conexiones de carga en vez de las \ConexionesTotales{} actuales. Por eso se recomienda un piloto en un sector, con holgura de energía ante atrasos y evaluando con RED la ampliación de la carga nocturna.
